\title{How to Use CLV3 Class File}

\begin{abstract}
This article describes how to use the ``CLV3'' class file, developed by {\sf SPi}, to produce typeset papers based on Computer Linguistics typesetting specifications for submission to MIT. This document was produced using ``{\tt clv3}'' class file and will be modified whenever there is an update in the layout specifications.
\end{abstract}

\section{Introduction}

This document applies to version 3 of CL class file. Prior style files such as 
``{\tt cl.sty}'' and ``{\tt coli.sty}'' do not have all of the features described here. It is assumed that the user has a basic knowledge of \LaTeX\ typesetting commands.

\section{Class File Options}

There are several options that can be used to switch the mode of MIT2 from normal article to manuscript style, or to different layout styles. This is specified in the usual \LaTeX\ way by declaring:

\verb|\documentclass[bookreview,manuscript]{clv3}|

\begin{deflist}
\item[bookreview] Sets the article layout for Book Review.
\item[brief] Sets the article layout for Briefly Noted.
\item[discussion] Sets the article layout for Squibs and Discussions.
\item[pubrec] Sets the article layout for Publication Received.
\item[shortpaper] Sets the article layout for Short Paper.
\item[manuscript] Sets the baseline spacing to double space. This option can be used in combination with other options.
\end{deflist}

By not declaring any option in the \verb|\documentclass| command the class file will automatically set to standard article layout.

\section{Title Page}

The title page is created using the standard \LaTeX\ command \verb|\maketitle|. 
Before this command is declared, the author must declare all the data which are to appear in the title area.\footnote{$\backslash$maketitle is the command to execute all the title page information.}

\subsection{Volume, Number and Year}

The command \verb|\issue{vv}{nn}{yyyy}| is used in declaring the volume, number and year of the article. The first argument is for the volume, the second argument is for the issue number. Volume and Issue number will appear on the even page running head opposite the journal name. The third argument is for the Year which will appear in the copyright line at the bottom of the title page.

\subsection{Document Head}

Document head is produced with the command \verb|\dochead{Document Head}|. Doc head will output differently, or may not appear at all, depending on the option used in the documentclass.

\subsection{Paper Title}

The paper title is declared like: \verb|\title{Computer Linguistic Article}| in the usual \LaTeX manner. Line breaks may be inserted with (\verb|\\|) to equalize the length of the title lines.

\subsection{Authors}
The name and related information for authors is declared with the \verb|\author{}| command. 

The \verb|\thanks{}| command produces the ``first footnotes.''. \LaTeX\ \verb|\thanks| 
cannot accommodate multiple paragraphs, author will have to use a separate \verb|\thanks|
for each paragraph.

The \verb|\affil{}| command produces the author affiliations that appears right under the author's name.

\subsection{Running Headers} The running heads are declared with the \verb|\runningtitle{Running Title}| for the journal name and \verb|\runningauthor{Author's Surname}| for author. These information will appear on the odd pages. For {\tt bookreview} option, odd page running head is automatically set to "Book Reviews". Even page running head is default to Computational Linguistics opposite volume and issue number.

\subsection{History Dates}

History dates are declared with \verb|\historydates{Submission received:...}|. This data should contain Submission, Revised and Accepted date of the article. History dates appear at the footnote area of title age.

\section{Abstract}

Abstract is the first part of a paper after \verb|\maketitle|. Abstract text is placed within the abstract environment:

\begin{verbatim}
\begin{abstract}
This is the abstract text . . .
\end{abstract}
\end{verbatim}

\section{Section Headings}

Section headings are declared in the usual \LaTeX\ way via \verb|\section{}|, 
\verb|\subsection{}|, \verb|\subsubsection{}|, and \verb|\paragraph{}|. The first 3 levels of section head will have Arabic numbering separated by period. The \verb|\paragraph{}| section will have the title head in Italics and at the same line with the first line of succeeding paragraph.

\section{Citations}
Citations in parentheses are declared using the \verb|\cite{}|
command, and appear in the text as follows: This technique is widely used \cite{woods}. The command \verb|\citep{}| (cite parenthetical) is a synonym of \verb|\cite{}|.

Citations used in the sentence are declared using the \verb|\namecite{}|
commands, and appear in the text as follows: 
\namecite{woods} first described this technique. The command \verb|\citet{}| (cite textual) is a synonym of \verb|\namecite{}|.

This style file is designed to be used with the BibTeX style file \verb|compling.bst|.  Include the command
\verb|
\end{document}
